\SourceLeaf{445}%
\[
(\xi\xi'+\eta\eta'+\zeta\zeta')\frac{\partial t}{\partial\theta}=\frac{\partial x}{\partial s}.\frac{\partial x}{\partial\theta}+\frac{\partial y}{\partial s}.\frac{\partial y}{\partial\theta}+\frac{\partial z}{\partial s}.\frac{\partial z}{\partial\theta},
\]
welche Grösse wir mit $u$ bezeichnen wollen, und die also ihrerseits wieder Function von $\theta$ und $s$ sein wird.

Man findet dann, wenn man nach $s$ differentiirt:
\[
\begin{aligned}
\frac{\partial u}{\partial s}={}&\frac{\partial\partial x}{\partial s^{2}}.\frac{\partial x}{\partial\theta}+\frac{\partial\partial y}{\partial s^{2}}.\frac{\partial y}{\partial\theta}+\frac{\partial\partial z}{\partial s^{2}}.\frac{\partial z}{\partial\theta}\\
&+\tfrac12\frac{\partial\left\{\left(\frac{\partial x}{\partial s}\right)^{2}+\left(\frac{\partial y}{\partial s}\right)^{2}+\left(\frac{\partial z}{\partial s}\right)^{2}\right\}}{\partial\theta}\\
={}&\frac{\partial\partial x}{\partial s^{2}}.\frac{\partial x}{\partial\theta}+\frac{\partial\partial y}{\partial s^{2}}.\frac{\partial y}{\partial\theta}+\frac{\partial\partial z}{\partial s^{2}}.\frac{\partial z}{\partial\theta},
\end{aligned}
\]
weil
\[
\left(\frac{\partial x}{\partial s}\right)^{2}+\left(\frac{\partial y}{\partial s}\right)^{2}+\left(\frac{\partial z}{\partial s}\right)^{2}=1,
\]
also sein Differential $=0$.

Allein da alle zu constantem $\theta$ [gehörigen] Punkte in einer kürzesten Linie liegen, so ist, wenn $L$ das Zenith des Punktes, welchem $s,\theta$ entsprechen, bedeutet und die Coordinaten von $L$ durch $X,Y,Z$ bezeichnet werden:
\[
\frac{\partial\partial x}{\partial s^{2}}=\frac{X}{p},\qquad\frac{\partial\partial y}{\partial s^{2}}=\frac{Y}{p},\qquad\frac{\partial\partial z}{\partial s^{2}}=\frac{Z}{p},
\]
wenn $p$ den Krümmungshalbmesser bedeutet. Wir haben also
\[
p.\frac{\partial u}{\partial s}=X\frac{\partial x}{\partial\theta}+Y\frac{\partial y}{\partial\theta}+Z\frac{\partial z}{\partial\theta}=\frac{\partial t}{\partial\theta}(X\xi'+Y\eta'+Z\zeta').
\]
Allein
\[
X\xi'+Y\eta'+Z\zeta'=\cos L\lambda'=0,
\]
weil offenbar $\lambda'$ in dem grössten Kreise liegt, dessen Pol $L$ ist. Wir haben daher
\[
\frac{\partial u}{\partial s}=0
\]
oder $u$ von $s$ unabhängig und daher bloss Function von $\theta$. Allein für $s=0$ ist offenbar $t=0$, $\frac{\partial t}{\partial\theta}=0$ und also $u=0$: wir schliessen daraus, dass allgemein $u=0$ sein wird oder
\[
\cos\lambda\lambda'=0.
\]
Es folgt hieraus der schöne Lehrsatz:
